% main.tex - the coordination-tax paper. Hand-written; gen-tex.mjs leaves it alone.
% Bibliography is generated: ../refs.bib comes from frontend/src/sources.ts.
\documentclass[conference]{IEEEtran}
\usepackage{graphicx}
\usepackage{cite}
\usepackage{url}
\usepackage{amsmath}

\title{The Coordination Tax: Light-Speed Delay in a Self-Replicating Probe Swarm Bites Only When Hops Are Long}

\author{
\IEEEauthorblockN{Noah Hyd\'en}
\IEEEauthorblockA{Independent researcher \\ ORCID: 0009-0003-4523-0467}
}

\begin{document}
\maketitle

\begin{abstract}
A swarm of self-replicating probes fills a stellar field star by star, each settling
a star and launching copies. The standard exploration-timescale model grants every
probe perfect, instantaneous knowledge of which stars are already settled; finite
light-speed is its stated future work. We add that light-speed limit: a probe deciding
at a star treats a distant star as settled only once the news, travelling at $c$, has
arrived, so probes race for the same target from stale views and waste trips. Over a
32-seed paired ensemble on identical galaxies, the delay slows the time to fill the
field by a median of about 0\% for powered nearest-neighbour flight, about 30\% for
nearest-star slingshot flight, and about 50\% for maximum-boost slingshot flight; every
case still reaches 100\%, so lag alone produces no permanent unsettled fraction. The
lag scale is set by the dimensionless ratio of information-delay to travel-time per
hop, which equals the probe speed in units of $c$; whether it actually bites is decided
by how non-local each hop is. Fast but local flight is nearly immune, because a probe
that loses a race simply takes the star next door; the cost falls on long-range hops
made from stale views. This refines the perfect-information picture: the slingshot
speed-up is real, but under finite light-speed a meaningful fraction of it is eaten by
uncoordinated long-range collisions.
\end{abstract}

\begin{IEEEkeywords}
self-replication, interstellar exploration, light-speed delay, multi-agent coordination, swarm dynamics
\end{IEEEkeywords}

\section{Introduction}

A self-replicating probe does not explore a galaxy alone. It settles a star, builds
copies of itself, and sends them on to further stars, so a single launch grows into a
front that sweeps outward and fills the reachable field. How long that takes is the
exploration-timescale question, and the sharpest treatment of it for self-replicators is
Nicholson and Forgan's slingshot study \cite{nicholson-forgan-2013}: probes that steal
energy from the motion of the stars they pass (a gravitational slingshot, following the
single-probe analysis of Forgan, Papadogiannakis and Kitching
\cite{forgan-papadog-kitching-2013}) can explore roughly two orders of magnitude faster
than probes under their own power.

That result rests on one convenient assumption, which its authors name as a limitation
and leave for future work: every probe is granted perfect, instantaneous knowledge of
which stars are already settled. No probe ever flies to a star another probe has just
claimed, because it always knows. In a real galaxy no such oracle exists. News that a
star has been settled can travel no faster than light, so a probe choosing its next
target is acting on a picture of the field that is, for distant stars, years or
centuries out of date.

This paper adds that light-speed limit to the model and measures what it costs. The
question is not merely academic bookkeeping: it decides whether the celebrated slingshot
speed-up survives contact with the finite speed of information, or whether a swarm racing
outward on stale maps wastes much of the advantage colliding with itself. The answer,
built from a controlled paired-ensemble experiment, is that the tax is real but sharply
selective - it is set in scale by the probe's speed and in bite by how far each hop
reaches.

\section{The light-delayed belief model}
\label{sec:model}

\subsection{The knowledge gate}
The change to the model is one gate on knowledge. A settled star is treated as an
omnidirectional beacon that begins emitting in the year it is settled. A probe deciding
its next target at star $f$ in year $Y$ believes a distant star $i$ is settled only once
that beacon's light has had time to arrive:
\begin{equation}
\text{settled}(i) + \frac{\mathrm{dist}(f, i)}{c} \le Y . \label{eq:gate}
\end{equation}
Under the perfect-information assumption of the source model, (\ref{eq:gate}) collapses to
``settled at any non-negative year,'' and the model reproduces our perfect-information
baseline bit-for-bit; the light-speed term is what introduces stale views. When two probes
both believe an unsettled star is the best target they both fly to it; the loser arrives to
find it taken, its trip wasted, and must re-target from the even more stale picture it holds
on arrival. This is availability chosen over consistency: rather than wait for a single
agreed map that finite light-speed makes impossible, each probe acts now on its local view
\cite{lamport-1978}. Stars carry velocity magnitudes that feed the slingshot boost but are
held fixed in position, following the source model, so $\mathrm{dist}(f,i)$ in
(\ref{eq:gate}) is static and the beacon geometry does not shift under the probe.

\subsection{The three flight policies}
We carry three named policies, following the slingshot studies we extend
\cite{forgan-papadog-kitching-2013,nicholson-forgan-2013}. \emph{Powered
nearest-neighbour} flight moves under the probe's own power at a cruise speed of
$3\times10^{-5}\,c$ to the nearest unsettled star. \emph{Nearest-star slingshot} flight
takes the same nearest target but gains speed from gravitational slingshots off the stars
it passes. \emph{Maximum-boost slingshot} flight instead selects, among the nearest
believed-unsettled stars, the target that yields the largest velocity gain, deliberately
reaching past near neighbours to distant, fast stars. The three span the axis that matters
here: powered flight is slow and local, nearest-star slingshot is fast and local, and
maximum-boost slingshot is fast and far. The effective speeds the policies reach in the
simulation, which set the lag ratio below and let the reader check the reported penalties,
are given with the results in Table~\ref{tab:mech}: a powered cruise near $9$ km/s against
slingshot means near $2{,}700$ km/s (nearest) and $3{,}900$ km/s (maximum-boost).

\subsection{The lag ratio, and what it can and cannot say}
How much the gate can matter on a single hop is governed by a dimensionless ratio: the
information delay across the hop divided by the travel time across the same hop. For a hop
of length $d$ the delay is $d/c$ and the travel time is $d/v$, so the ratio is
\begin{equation}
\Lambda \;=\; \frac{d/c}{d/v} \;=\; \frac{v}{c}, \label{eq:lambda}
\end{equation}
the probe's speed in units of $c$, and it is the \emph{same on every hop}: the length $d$
cancels. A slow probe ($\Lambda \ll 1$) outpaces its own ignorance, so stale views barely
arise; at the powered cruise speed $\Lambda \approx 3\times10^{-5}$ is negligible. A fast
probe carries a large $\Lambda$ and routinely decides before the relevant news can reach it.
So $\Lambda$ gates \emph{whether} the delay can bite at all, and it does so correctly: the
effect appears only once probes are fast, which is the slingshot regime. But precisely
because $\Lambda$ is hop-length-independent it cannot say \emph{where} the cost lands, and
Section~\ref{sec:results} shows the penalty is decided by hop length, not by $\Lambda$. We
therefore keep $\Lambda$ for its gating role and demote it as a predictor of the penalty's
size.

\subsection{Experimental design and parameters}
The field is a box of stars at a uniform density of one star per cubic parsec, the
convention of the source model. Each settled star launches two offspring (the replication
branching factor); this is what makes the field fill exponentially rather than as a single
roving chain, and it sets the swarm density that governs how often two probes contend for
one star. Three properties keep the measurement honest. First, knowledge is evaluated only
at the decision star at the moment of decision; news a probe's path happens to cross in
mid-flight is ignored. This deliberately undercounts what a probe could know, so the measured
waste is a conservative upper bound on the true penalty. Second, a probe that loses a race
re-targets, but only up to a fixed number of times before it is retired as wasted; this cap
bounds pathological bounce chains late in the fill, is bookkeeping rather than physics, and
the results are insensitive to it. Third, the whole simulation is a pure, seeded, fixed-step
fold: fixing the seed fixes the star field, every arrival, and every target choice, so a run
with the light-speed gate and a run without it can be compared on the \emph{identical} galaxy.
That paired design is what lets us attribute a difference in fill-time to the delay alone
rather than to the luck of the star layout.

The uniform density is easy to over-read, so one clarification. Changing the \emph{value} of
the uniform density rescales every inter-star distance by a common factor, which multiplies
both the travel time and the light-travel time on every hop equally; $\Lambda$ and, in the
continuum limit, the relative penalty are left unchanged (a pure geometric rescaling of the
absolute clock). We confirm this holds in practice: re-running at densities down to the
sparser solar-neighbourhood value near $0.14$ stars per cubic parsec \cite{recons-census}
keeps the median penalty within a few points of its value at unit density (about
\mbox{27-29\%} for nearest-star, about \mbox{46-52\%} for maximum-boost), the residual scatter
being the fixed timestep, not the density. A \emph{non-uniform}, clumpier distribution is a different matter, and we
treat it as a limitation (Section~\ref{sec:limitations}): by lengthening the long-range hops
that carry the penalty it is expected to move the relative penalties, not merely the absolute
timescale.

\subsection{Baseline validation}
The paired comparison is only as good as its perfect-information baseline, so we check that
baseline at two levels. By construction the ``instant'' regime is the $c\to\infty$ limit of
(\ref{eq:gate}), and it reproduces the default perfect-information fold bit-for-bit. Against
Nicholson and Forgan's published model \cite{nicholson-forgan-2013} we reproduce their two
qualitative findings: gravitational slingshots explore far faster than powered flight, and
their surprising result that aiming for the nearest star beats aiming for the largest boost
on total exploration time. In our fold, on a $400$-star field, nearest-star slingshot fills
in about $75{,}000$ years against about $305{,}000$ for maximum-boost, while maximum-boost
reaches the higher peak speed. We do not reproduce their roughly hundred-fold slingshot
speed-up quantitatively: at the fixed timestep the boosted hops are shorter than one step and
are quantized, so the measured speed-up is about twenty-fold. Lowering the timestep recovers
more of it; the qualitative ordering, on which this paper's conclusions rest, is
timestep-robust.

\section{Results: the coordination tax}
\label{sec:results}

We measured the cost as a paired ensemble: 32 seeded galaxies of 300 stars each, every one
filled twice, once with perfect information and once with the light-speed gate of
(\ref{eq:gate}), so that the only difference between the two runs of a pair is the delay.
The primary quantity is the relative slowdown in the time to settle the entire field,
$t_{100}$, per seed; because $t_{100}$ is a tail statistic dominated by the last few stars,
we also report the penalty at earlier coverage fractions and across a range of field sizes.
Reporting the distribution across the ensemble rather than a single run follows the seeded
Monte Carlo practice of the field \cite{forgan-2009}.

\begin{table}[!t]
  \centering
  \caption{Fill-100\% penalty from light-speed lag over the 32-seed paired ensemble at
    300 stars: median, interquartile range, and a seeded percentile-bootstrap 95\%
    confidence interval on the median. The last column is a two-sided sign test over the
    paired seeds; the slingshot penalty is positive in all 32, while powered is null
    (positive in only 6 of the 11 seeds that differ at all).}
  \label{tab:penalty}
  \begin{tabular}{lcccc}
    \toprule
    Flight policy & Median (\%) & IQR (\%) & 95\% CI (\%) & $p$ \\
    \midrule
    Powered nearest-neighbour & 0.0  & [0.0, 0.0]   & [0.0, 0.0]   & 1.0 \\
    Nearest-star slingshot    & 30.3 & [20.0, 38.5] & [23.5, 37.5] & $<10^{-9}$ \\
    Maximum-boost slingshot   & 51.4 & [46.0, 54.0] & [47.2, 53.7] & $<10^{-9}$ \\
    \bottomrule
  \end{tabular}
\end{table}

\begin{table}[!t]
  \centering
  \caption{The mechanism, as medians over the 32 light-speed runs. Effective speed is the
    mean speed at which probes are launched; $\Lambda_{\text{eff}} = v_{\text{eff}}/c$. Hop
    length is the mean distance of a wasted (lost-race) trip and of a winning (settling)
    trip. The two slingshot policies differ in $\Lambda_{\text{eff}}$ by about
    $1.4\times$ but in wasted-hop length by about $1.85\times$, and the penalty in
    Table~\ref{tab:penalty} tracks the latter.}
  \label{tab:mech}
  \begin{tabular}{lcccc}
    \toprule
    Flight policy & $v_{\text{eff}}$ (km/s) & $\Lambda_{\text{eff}}$ & Wasted (pc) & Settled (pc) \\
    \midrule
    Powered nearest-neighbour & 9     & $3.0\times10^{-5}$ & 1.62 & 0.83 \\
    Nearest-star slingshot    & 2{,}734 & $9.1\times10^{-3}$ & 1.44 & 0.83 \\
    Maximum-boost slingshot   & 3{,}897 & $1.3\times10^{-2}$ & 2.67 & 2.36 \\
    \bottomrule
  \end{tabular}
\end{table}

The medians and their spread are given in Table~\ref{tab:penalty} and shown per seed in
Fig.~\ref{fig:slowdown}. Powered flight pays essentially nothing; nearest-star slingshot
loses about a third of its time; maximum-boost slingshot loses about half. The two slingshot
penalties are not marginal: each is positive in all 32 seeds, and the bootstrap confidence
intervals on the medians (Table~\ref{tab:penalty}) are well clear of zero and of each other.

The central finding is that the penalty is \emph{not} governed by $\Lambda = v/c$ from
(\ref{eq:lambda}). Table~\ref{tab:mech} makes the case directly. The two slingshot policies
reach comparable effective speeds, $2{,}734$ against $3{,}897$ km/s, so their
$\Lambda_{\text{eff}}$ differ by only about $1.4\times$; yet their penalties differ by about
$1.7\times$. If $\Lambda$ set the cost, the faster policy would pay in that same $1.4$-to-$1$
proportion, and the powered policy, though slowest, would still pay something. Neither holds.
What does track the penalty is hop non-locality. A probe flying to the \emph{nearest} star
that loses a race has lost almost nothing: the next unsettled star is right beside it, so the
wasted trip and the recovery are both short. A probe chasing the \emph{biggest boost}
deliberately reaches past its neighbours to a distant star, so when it loses that race both
the wasted trip and the re-targeting detour are long. The simulation confirms the mechanism
in the hop lengths themselves: the mean wasted trip is $2.67$ pc for maximum-boost against
$1.44$ pc for nearest-star, and its winning trips are nearly three times longer too
(Table~\ref{tab:mech}). So $\Lambda$ decides only \emph{whether} lag can bite: powered flight,
with $\Lambda \approx 3\times10^{-5}$, is immune not because its wasted hops are short (they
are not, at $1.62$ pc) but because it is far too slow for any race to matter. Given that a
policy is fast enough for lag to bite, the \emph{non-locality of its hops} sets how hard.
Fig.~\ref{fig:settlement} shows the effect on a single representative galaxy for the
nearest-star policy: the light-speed run reaches full settlement at about 110{,}000 years
against 80{,}000 for the perfect-information run on the same seed, a right-shift of roughly a
third visible as a lag in the settlement curve, not a change in its final height.

The penalty is not an artifact of the fragile $t_{100}$ tail. Fig.~\ref{fig:coverage} reports
the median penalty at coverage fractions from $t_{25}$ to $t_{100}$: for nearest-star flight
it rises steadily from about 14\% at quarter-coverage to about 30\% at full coverage, and for
maximum-boost it is already about 26\% at quarter-coverage and stays near or above 50\%
thereafter. The shift is present throughout the fill, and grows as the front reaches farther
from the origin, exactly as the locality mechanism predicts; it is not a single unlucky last
star. Powered flight stays at zero at every fraction.

Nor is it an artifact of the small $300$-star box, the concern that matters most for
extrapolating toward galactic saturation. Fig.~\ref{fig:finitesize} sweeps the system size
across a four-fold range: the powered penalty stays near zero, the nearest-star penalty stays
near a third, and the maximum-boost penalty holds near or slightly above a half, firming up
rather than washing out as the field grows and hops lengthen. The maximum-boost policy scans
its nearest neighbours at every launch and so costs $O(N^2)$, which bounds the range we sweep;
within that range the penalties are stable, and the trend is toward more penalty at larger
scale, not less.

Finally, what does \emph{not} happen: every case still fills the field to full coverage. This
is a property of the model, not a result to be discovered. Re-targeting guarantees that a
probe which loses a race simply picks another star, and with no loss term the reachable field
must eventually close; light-speed lag makes the swarm slower and more wasteful but leaves no
permanent unsettled remainder. It is worth stating precisely because it marks the boundary of
what delay alone can do: a permanent settled fraction below one, an ``Aurora'' steady state in
which some region stays forever empty \cite{carroll-nellenback-2019}, would require settled
sites to \emph{die}, and delay supplies no such term. Lag is a tax on speed, not a barrier to
completion.

\begin{figure}[!t]
  \centering
  \includegraphics[width=\columnwidth]{fig_slowdown_by_policy.pdf}
  \caption{Distribution of the per-seed fill-100\% penalty imposed by light-speed-limited
    coordination, over the fixed 32-seed ensemble at 300 stars, for the three flight policies
    (box: interquartile range and median; points: all 32 individual seeds). Powered
    nearest-neighbour flight pays essentially no penalty, because a lost race is recovered at
    the star next door, while the long-range maximum-boost policy fills the field materially
    later, because wasted trips launched from stale views become long detours.}
  \label{fig:slowdown}
\end{figure}

\begin{figure}[!t]
  \centering
  \includegraphics[width=\columnwidth]{fig_penalty_by_coverage.pdf}
  \caption{Median fill-time penalty at coverage fractions from $t_{25}$ to $t_{100}$, over the
    32-seed ensemble, for the three policies. The slowdown is present across the whole fill,
    not only in the $t_{100}$ tail, and grows as the front reaches farther from the origin;
    powered flight stays at zero throughout.}
  \label{fig:coverage}
\end{figure}

\begin{figure}[!t]
  \centering
  \includegraphics[width=\columnwidth]{fig_finite_size.pdf}
  \caption{Median fill-100\% penalty versus system size for the three policies (line: median;
    shaded band: interquartile range over 8 seeds per point). Across a four-fold increase in
    the number of stars the penalties are stable, holding near 0\%, a third, and a half
    respectively, and firm up rather than wash out. The $O(N^2)$ cost of the maximum-boost
    scan bounds the range swept.}
  \label{fig:finitesize}
\end{figure}

\begin{figure}[!t]
  \centering
  \includegraphics[width=\columnwidth]{fig_settlement_curves.pdf}
  \caption{Fraction of the reachable field settled versus year for the nearest-star slingshot
    policy on a single representative galaxy, under perfect global information (solid) and
    under light-speed-limited coordination (dashed). The light-speed curve is shifted later,
    reaching full settlement at about 110{,}000 years versus 80{,}000, a delay attributable to
    information lag alone; both curves reach the same final height.}
  \label{fig:settlement}
\end{figure}

\section{Discussion}

The result refines rather than overturns the slingshot picture. Nicholson and Forgan's
speed-up is real and large \cite{nicholson-forgan-2013}, but it was computed under perfect
coordination, and coordination is not free. Under finite light-speed a meaningful fraction
of the advantage, roughly a third for nearest-star slingshots and half for maximum-boost,
is spent on trips to stars that turn out to be already taken. The more aggressively a
policy reaches for distant, high-boost targets, the more of its speed it loses to
uncoordinated collisions.

\subsection{Limitations}
The measurement is a conservative upper bound in two senses, and both point to future work.
Probes here learn only at their decision stars, ignoring news their paths cross in flight,
so a two-endpoint knowledge model would recover some of the wasted trips; and probes here
share nothing directly, so adding probe-to-probe relay of the settled map would let a swarm
partially coordinate rather than merely collide. Both would lower the tax without erasing
it, since the fundamental limit, that no probe can know a distant star is taken until its
light arrives, is set by physics, not by protocol. Two further simplifications bound the
scope. The field is at a uniform one star per cubic parsec; a realistic, clumpier
distribution would lengthen the long-range hops that carry the penalty, so the reported
figures are conservative for the maximum-boost policy in particular. And the dynamics have
no loss term: settled sites never fail, so the field always closes. A settlement-death term
of the kind that produces the ``Aurora'' steady state, in which the settled fraction
equilibrates below one \cite{carroll-nellenback-2019}, is exactly what would be needed to
turn lag into a permanent unsettled remainder; delay alone does not.

\subsection{The signature of delayed agreement}
The behaviour is the expected signature of delayed multi-agent agreement. Formal consensus
degrades as the ratio of communication delay to decision cadence grows: the standard
protocol is stable only while the one-hop delay stays below a bound set by the coupling
\cite{olfati-saber-murray-2004}, a result later generalized to switching topologies and
link failures \cite{olfati-saber-fax-murray-2007}. Our probes never reach agreement about
the settled map at all; at a mean inter-star spacing of about a parsec the round-trip
light-time is several years, far longer than any decision cadence, so they behave like
agents forced to act on purely local, out-of-date views. In the limit of unbounded delay,
guaranteed agreement is not merely slow but impossible \cite{flp-1985}, and a system that
must nonetheless keep acting is forced to trade a single consistent map for availability
\cite{gilbert-lynch-2002}, which is precisely the choice (\ref{eq:gate}) encodes.

This is the same progression that drove the space teleoperation and networking communities
away from real-time control. Past about a second of round-trip delay a human operator
abandons closed-loop control for move-and-wait \cite{ferrell-1965}, and as delay grows the
handoff continues through supervisory control to full remote autonomy \cite{sheridan-1993}.
Past minutes to hours no continuous end-to-end path exists at all, and communication becomes
hop-by-hop store-and-forward, the delay-tolerant regime motivated by the interplanetary case
\cite{burleigh-2003} and standardized first as an architecture \cite{rfc-4838} and now as a
deployed protocol \cite{rfc-9171}. At interstellar hop lengths every one of those faster
regimes is left far behind; each probe is effectively an independent colony acting on
priors, and the coordination tax is what that independence costs when probes want the same
stars.

\subsection{A design lesson}
There is a practical corollary. For a real swarm on stale maps, a locally greedy policy is
not merely simpler than a boost-maximising one; it is more robust to the delay it cannot
avoid, because a lost race costs only the star next door. A designer who cannot buy
coordination should prefer locality, spending the slingshot advantage on speed rather than
on reach.


\bibliographystyle{IEEEtran}
\bibliography{../refs}

\end{document}
